\chapter{The even--even sector}
\label{chap:even-even}

Even reflection across the separatrix gives a zero normal trace.
The normal-field virial then has nonnegative bulk density, while transport
and a one-dimensional sum of squares make its limiting boundary term
nonpositive. For the exclusion, fix $\lambda\ne0$ and first work in one
rotation character $\kappa=\pm1$; Lemma~\ref{lem:rotation-characters}
reduces the sector to these two cases.

\emph{Proof roadmap.}
\begin{enumerate}
  \item \textbf{Truncated virial}
        (Section~\ref{sec:ee-virial}).
        Apply the $L^2$ virial on annuli and remove the center boundary to
        obtain $\int_{\{\Psi>2h\}}\mathcal J_\omega=G_\omega(h)$
        for almost every $h$, with $\mathcal J_\omega\ge0$.
  \item \textbf{Boundary limits}
        (Section~\ref{sec:ee-gradient-limits}).
        Neumann parity kills the normal energy; the tangential energy
        converges to the weighted arc energy. Transport supplies a common
        sequence $h_j\to0$, avoiding all exceptional levels, on which
        $G_\omega(h_j)\to-2\sum_k\mathcal A_\lambda(g_k)$
        (Lemma~\ref{lem:ee-transport-traces}).
  \item \textbf{Boundary positivity}
        (Section~\ref{sec:ee-boundary-positivity}).
        The canonical arc unknown lies in the weighted interval domain.
        Completion of squares and affine splitting give
        $\mathcal A_\lambda(g_k)\ge0$: the endpoint relation from
        $\kappa=\pm1$ cancels the indefinite affine cross term.
        No vanishing saddle values are assumed.
  \item \textbf{Vanishing}
        (Section~\ref{sec:ee-assembly}).
        Nonnegative exhaustion forces $\mathcal J_\omega=0$.
        Since $f>0$ almost everywhere in the cell, $\omega=0$ there;
        reflection gives vanishing on the torus
        (Lemma~\ref{lem:ee-character-exclusion}).
\end{enumerate}

The interval boundary calculus is reusable independently of this sector.
The gradient traces and their parity are sector-specific; the lift and
transport interfaces state their own hypotheses. Only truncated bulk
integrability is used, not integrability up to the saddles.

\section{The normal field and truncated virial}
\label{sec:ee-virial}

\begin{definition}[Normal field]
\label{def:normal-field}
\uses{def:cells, lem:dipole-critical-points, lem:dipole-identities}
\lean{EulerDipoleEigenvalues.normalField₁,
  EulerDipoleEigenvalues.normalField₂,
  EulerDipoleEigenvalues.normalCoeff,
  EulerDipoleEigenvalues.normalC}
\leanok
On $\Cell$ minus its center, with $\rho^2=|\nabla\Psi|^2$ and
$\delta=\cos x-\cos y$, let (manuscript \texttt{(ee:normal-field)})
\[
  W=(\Psi-2)\frac{\nabla\Psi}{\rho^2},\qquad
  f=\frac{\Psi(2-\Psi)\delta^2}{\rho^4},\qquad
  c=\frac{\Psi(2-\Psi)}{\rho^2}.
\]
\end{definition}

\begin{lemma}[Normal field algebra]
\label{lem:normal-field-algebra}
\uses{def:normal-field, def:adapted-field}
\lean{EulerDipoleEigenvalues.isAdaptedOn_normalField,
  EulerDipoleEigenvalues.div_normalField,
  EulerDipoleEigenvalues.normalJacForm,
  EulerDipoleEigenvalues.normalJacForm_normal,
  EulerDipoleEigenvalues.normalJacForm_tangent,
  EulerDipoleEigenvalues.normalJacForm_mixed,
  EulerDipoleEigenvalues.gradSq_dipoleLift_pos,
  EulerDipoleEigenvalues.normalC_pos_le_one}
\leanok
$W$ is adapted with coefficient $f$, $\Div W=f+1$, and in the orthonormal
normal--tangent frame $(\nabla\Psi/\rho,\nabla^\perp\Psi/\rho)$,
\[
  2\sym DW-fI=\begin{pmatrix}2-c&0\\0&c\end{pmatrix}.
\]
Moreover $0<c\le1$ on $\Cell$ minus its center (manuscript
\texttt{(ee:normal-algebra)}, \texttt{(ee:c-positive)}).
\end{lemma}

\begin{proof}
\uses{lem:ee-arc-energy, lem:ee-transport-traces, lem:arc-trace-ac,
 lem:canonical-boundary-regularity, lem:separatrix-arc,
 lem:continuous-representative, lem:rotation-characters}
\leanok
The rotation character gives $g=\kappa^k\,\tilde\psi\circ\gamma$, so the
arc absolute-continuity, Melnikov and weighted-energy results supply the
hypotheses of Lemma~\ref{lem:canonical-boundary-regularity}.
Continuity on the compact torus supplies boundedness and the saddle limits.
The lemma gives the weighted domain, forcing regularity and endpoint values.
Finally $R(\pi,0)=(0,-\pi)$ and the character of $\tilde\psi$ give
$c_+=\kappa c_-$.
\end{proof}

\begin{remark}
Lean uses the unnormalized frame $(\nabla\Psi,\nabla^\perp\Psi)$, whose
vectors have squared length $\rho^2$, so its frame identities carry an
extra factor $\rho^2$. On the punctured cell,
$\rho^2\ge\Psi(2-\Psi)>0$.
\end{remark}

\begin{remark}[Digest note: regularity at the center]
The manuscript applies the virial lemma to fields reaching the center,
requiring $W\in W^{3,\infty}$ and $f\in W^{2,\infty}$ there. A symbolic
local expansion gives
$f=(x^2-y^2)^2/(4(x^2+y^2))+O((x^2+y^2)^2)$, whose nonpolynomial
leading term warns against assuming smoothness.
The proof avoids this issue: apply the virial only on annuli
(Lemma~\ref{lem:ee-annular-virial}), then remove the inner boundary by
Lemma~\ref{lem:ee-center-limit}. No higher regularity at the center is used.
\end{remark}

\begin{lemma}[The normal field cut off to an annulus]
\label{lem:normal-field-annulus}
\uses{def:normal-field, def:adapted-field}
\lean{EulerDipoleEigenvalues.annularCutoff_normalField}
\leanok
Let $0<\epsilon_1\le\epsilon_2<2$ and let $\chi$ be smooth and zero outside
$[\epsilon_1,\epsilon_2]$. Then $\chi(\Psi)W$ and $\chi(\Psi)f$ are smooth and periodic on the
lift, and $\chi(\Psi)W$ is adapted everywhere with coefficient $\chi(\Psi)f$.
\end{lemma}

\begin{proof}
\uses{lem:normal-field-algebra}
\leanok
On $\{0<\Psi<2\}$, $\rho^2\ge\Psi(2-\Psi)>0$, so $W$ and $f$ are smooth there; outside the
closed set $\{\epsilon_1\le\Psi\le\epsilon_2\}$ the products vanish identically near each point.
Adaptedness: $L_0(\chi(\Psi))=\chi'(\Psi)L_0\Psi=0$, so multiplying an adapted field by
$\chi(\Psi)$ multiplies its coefficient by $\chi(\Psi)$.
\end{proof}

\begin{lemma}[Cutoff algebra]
\label{lem:normal-field-cutoff-algebra}
\uses{def:normal-field, def:adapted-field, def:half-density-generator}
\lean{EulerDipoleEigenvalues.cutoff_normalField_algebra}
\leanok
In the setting of Lemma~\ref{lem:normal-field-annulus}, with $W_\chi=\chi(\Psi)W$ and
$f_\chi=\chi(\Psi)f$: $f_\chi-\Div W_\chi=-\chi(\Psi)-(\Psi-2)\chi'(\Psi)$, and for
$v\in\C^2$ with normal and tangential components $v_n,v_t$,
\[
  (2\sym DW_\chi-f_\chi I)v\cdot\bar v
  =\chi(\Psi)\bigl((2-c)|v_n|^2+c|v_t|^2\bigr)+2\chi'(\Psi)(\Psi-2)|v_n|^2 .
\]
\end{lemma}

\begin{proof}
\uses{lem:normal-field-algebra}
\leanok
$\Div(\chi W)=\chi\Div W+\chi'\nabla\Psi\cdot W=\chi(f+1)+\chi'(\Psi-2)$, and
$D(\chi W)=\chi DW+\chi'\,W\otimes\nabla\Psi$ with $W\otimes\nabla\Psi=(\Psi-2)\rho^{-2}\nabla\Psi\otimes\nabla\Psi$
symmetric; then use Lemma~\ref{lem:normal-field-algebra}. Both identities were checked
numerically at random points for a non-polynomial $\chi$ (October 2026).
\end{proof}

\begin{lemma}[Positivity of the bulk density]
\label{lem:ee-density}
\uses{lem:normal-field-algebra}
\lean{EulerDipoleEigenvalues.eeDensity,
  EulerDipoleEigenvalues.normalCoeff_nonneg,
  EulerDipoleEigenvalues.eeDensity_nonneg,
  EulerDipoleEigenvalues.eeDensity_gradient_pos}
\leanok
The density $\mathcal J=f|\omega|^2+(2-c)|\psi_n|^2+c|\psi_t|^2$ is
nonnegative on $\Cell$ minus its center, and its gradient part is strictly positive at every
regular interior point where $\nabla\psi\ne0$ (manuscript
\texttt{(ee:positive-density)}).
\end{lemma}

\begin{proof}
\uses{lem:normal-field-algebra}
\leanok
$f\ge0$ and $0<c\le1$.
\end{proof}

\begin{lemma}[A bound for the vorticity coefficient]
\label{lem:normal-coeff-bound}
\uses{def:normal-field}
\lean{EulerDipoleEigenvalues.normalCoeff_le}
\leanok
On $\Cell$ minus its center, $f\le(2-\Psi)/\Psi$. In particular $f$ is bounded on
$\{\Psi\ge\epsilon\}$ for every $\epsilon>0$, and $f\to0$ at the center.
\end{lemma}

\begin{proof}
\uses{lem:normal-field-algebra}
\leanok
Since $\Psi^2+\delta^2=2(\cos^2x+\cos^2y)=4-2\rho^2$, the claim
$\Psi(2-\Psi)\delta^2/\rho^4\le(2-\Psi)/\Psi$ is $\Psi^2(4-2\rho^2-\Psi^2)\le\rho^4$, that is
$(\rho^2+\Psi^2)^2\ge4\Psi^2$, which is $\rho^2\ge\Psi(2-\Psi)$
(Lemma~\ref{lem:normal-field-algebra}).
\end{proof}

\begin{definition}[Bulk density on an $L^2$ vorticity]
\label{def:ee-bulk-density}
\uses{lem:ee-density, def:stream-lifts, def:normal-tangent}
\lean{EulerDipoleEigenvalues.eeBulkDensity}
\leanok
For $\omega\in L^2(\T^2)$, write
$\mathcal J_\omega(p)=\mathcal J(p,\omega(p),\psi_n(p),\psi_t(p))$
on the lift, using the representatives of Definition~\ref{def:stream-lifts}.
\end{definition}

\begin{definition}[Level terms]
\label{def:ee-level-term}
\uses{def:level-integral, def:normal-tangent, def:stream-lifts}
\lean{EulerDipoleEigenvalues.eeLevelTerm}
\leanok
For $\omega\in L^2$ and $0<h<1$, with $\psi_n$ along $n=-\nabla\Psi/\rho$,
\[
  G_\omega(h)=\oint_{\{\Psi=2h\}}\operatorname{Re}(\psi\overline{\psi_n})\,d\ell
   +(2-2h)\left(\oint|\psi_n|^2\,dt-\oint|\psi_t|^2\,dt
     -\oint\operatorname{Re}(\omega\bar\psi)\,dt\right).
\]
\end{definition}

\begin{lemma}[Integrability of the bulk density]
\label{lem:ee-bulk-integrable}
\uses{def:ee-bulk-density}
\lean{EulerDipoleEigenvalues.integrableOn_eeBulkDensity,
  EulerDipoleEigenvalues.norm_normalDeriv_le,
  EulerDipoleEigenvalues.norm_tangentDeriv_le}
\leanok
For $\omega\in L^2$ and $h>0$, $\mathcal J_\omega$ is integrable on $\{p\in\Cell:\Psi>2h\}$.
Pointwise, $|\psi_n|,|\psi_t|\le|\psi_x|+|\psi_y|$.
\end{lemma}

\begin{proof}
\uses{lem:normal-coeff-bound, lem:normal-field-algebra, lem:lift-representatives}
\leanok
There $0\le f\le(2-\Psi)/\Psi\le1/h$ (Lemma~\ref{lem:normal-coeff-bound}) and $0<c\le1$; at
the center $f=c=0$. Since $|\sin x|,|\sin y|\le\rho$, $|\psi_n|,|\psi_t|\le|\psi_x|+|\psi_y|$, so
$\mathcal J_\omega\le|\omega|^2/h+6(|\psi_x|^2+|\psi_y|^2)$. Each square is integrable on the cell
by Lemma~\ref{lem:lift-representatives}, since $\omega,\psi_x,\psi_y\in L^2(\T^2)$.
\end{proof}

\begin{definition}[Boundary density]
\label{def:ee-boundary-density}
\uses{def:stream-lifts, def:normal-tangent}
\lean{EulerDipoleEigenvalues.eeBoundaryDensity}
\leanok
For $\omega\in L^2$, on the lift,
$E_\omega=\rho\operatorname{Re}(\psi\overline{\psi_n})+(2-\Psi)\bigl(|\psi_n|^2-|\psi_t|^2
-\operatorname{Re}(\omega\bar\psi)\bigr)$.
\end{definition}

\begin{lemma}[The boundary density]
\label{lem:ee-boundary-integrable}
\uses{def:ee-boundary-density, def:ee-level-term}
\lean{EulerDipoleEigenvalues.integrableOn_eeBoundaryDensity,
  EulerDipoleEigenvalues.ae_levelIntegral_eeBoundaryDensity}
\leanok
For $\omega\in L^2$, $E_\omega$ is integrable on the cell, and
$\oint_{\{\Psi=2h\}}E_\omega\,dt=G_\omega(h)$ for almost every $h\in(0,1)$.
\end{lemma}

\begin{proof}
\uses{lem:ee-bulk-integrable, lem:lift-representatives, lem:orbit-coordinates-area,
  lem:level-integral-congr}
\leanok
$|E_\omega|\le2\rho|\psi|(|\psi_x|+|\psi_y|)+2\bigl(2(|\psi_x|+|\psi_y|)^2+|\omega||\psi|\bigr)$, a sum of
products of functions in $L^2$ of the cell. Each of the four pieces
$\rho\operatorname{Re}(\psi\overline{\psi_n})$, $|\psi_n|^2$, $|\psi_t|^2$, $\operatorname{Re}(\omega\bar\psi)$
is dominated by a multiple of $|\psi|^2+|\psi_x|^2+|\psi_y|^2+|\omega|^2$, hence integrable on
the cell, and by Lemma~\ref{lem:orbit-coordinates-area} its restriction to almost every level is
integrable in flow time. On the level $\Psi=2h$ the factor $2-\Psi$ is $2-2h$, so the level
integral splits (Lemma~\ref{lem:level-integral-congr} for the period interval).
\end{proof}

\begin{lemma}[The cutoff identity]
\label{lem:ee-cutoff-identity}
\uses{def:ee-bulk-density, def:ee-boundary-density, def:eigen-problem}
\lean{EulerDipoleEigenvalues.ee_cutoff_identity}
\leanok
Let $\lambda\ne0$ and let $\chi$ be smooth and vanish outside a compact subinterval of
$(0,2)$. Then $\int_\Cell\chi(\Psi)\mathcal J_\omega=\int_\Cell\chi'(\Psi)E_\omega$.
\end{lemma}

\begin{proof}
\uses{lem:virial-l2, lem:normal-field-annulus, lem:normal-field-cutoff-algebra,
  lem:stream-ibp, lem:cell-periodic-integral, lem:lift-torus, lem:ee-bulk-integrable,
  lem:ee-boundary-integrable}
\leanok
Apply Lemma~\ref{lem:virial-l2} to $\chi(\Psi)W$ (Lemma~\ref{lem:normal-field-annulus}). By
Lemma~\ref{lem:normal-field-cutoff-algebra} its quadratic form is
\[
  0=\int\chi\mathcal J+2\int\chi'(\Psi-2)|\psi_n|^2
    +\tfrac12\int\Delta\bigl(-\chi-(\Psi-2)\chi'\bigr)|\psi|^2 .
\]
By Lemma~\ref{lem:stream-ibp} the last term is $-\int\chi K-\int(\Psi-2)\chi'K$, and
$\int\chi K=-\int\operatorname{Re}(\bar\psi\,\nabla\chi(\Psi)\cdot\nabla\psi)
=\int\chi'(\Psi)\rho\operatorname{Re}(\psi\overline{\psi_n})$ since $\nabla\Psi\cdot\nabla\psi=-\rho\psi_n$.
Where $\chi'(\Psi)\ne0$, $\rho>0$ and $|\nabla\psi|^2=|\psi_n|^2+|\psi_t|^2$, so
$2|\psi_n|^2-K=|\psi_n|^2-|\psi_t|^2-\operatorname{Re}(\omega\bar\psi)$. Collecting terms gives the
identity on the period square; all integrands are periodic and vanish where $\Psi\le0$, so
Lemma~\ref{lem:cell-periodic-integral} moves it to the cell.

In Lean the computation is organized pointwise: $\chi(\Psi)\mathcal J_\omega-\chi'(\Psi)E_\omega$
equals a weighted sum of ten integrands of the form $c\operatorname{Re}(u\bar v)$, $\phi K$
and $\operatorname{Re}(\bar\psi\,\nabla\chi(\Psi)\cdot\nabla\psi)$, each integrable on the
period square by Lemma~\ref{lem:lift-torus}. The virial identity and the two cases of
Lemma~\ref{lem:stream-ibp} make the weighted sum of their integrals vanish. Where
$\chi'(\Psi)\ne0$ the frame identities $\rho\operatorname{Re}(\psi\overline{\psi_n})
=\operatorname{Re}(\psi\,\overline{\sin x\,\psi_x+\sin y\,\psi_y})$ and
$|\psi_n|^2+|\psi_t|^2=|\nabla\psi|^2$ use $\rho>0$. On the cell both sides are integrable by
Lemmas~\ref{lem:ee-bulk-integrable} and~\ref{lem:ee-boundary-integrable}.
\end{proof}

\begin{lemma}[Virial identity on an annulus]
\label{lem:ee-annular-virial}
\uses{def:ee-bulk-density, def:ee-level-term, def:eigen-problem}
\lean{EulerDipoleEigenvalues.ee_annular_virial}
\leanok
Let $\lambda\ne0$. There is a null set $N$ such that for all $0<h<h'<1$ outside $N$,
\[
  \int_{\{2h<\Psi<2h'\}\cap\Cell}\mathcal J_\omega=G_\omega(h)-G_\omega(h').
\]
\end{lemma}

\begin{proof}
\uses{lem:ee-cutoff-identity, lem:ee-bulk-integrable, lem:ee-boundary-integrable,
  lem:coarea-level-cutoff, lem:du-bois-reymond}
\leanok
Write $\chi(s)=\varphi(s/2)$ with $\varphi$ smooth and vanishing outside $[a,b]\subset(0,1)$.
By Lemma~\ref{lem:coarea-level-cutoff}, the cutoff identity
(Lemma~\ref{lem:ee-cutoff-identity}) becomes $\int\varphi\,m=\int\varphi'\tilde G$ with
$m(h)=2\oint_{\{\Psi=2h\}}\mathcal J_\omega\,dt$ and $\tilde G(h)=\oint_{\{\Psi=2h\}}E_\omega\,dt$, both
integrable on compact subintervals (Lemmas~\ref{lem:ee-bulk-integrable}
and~\ref{lem:ee-boundary-integrable}). By Lemma~\ref{lem:du-bois-reymond},
$\tilde G(h)-\tilde G(h')=\int_h^{h'}m$ outside a null set, and the right side is
$\int_{\{2h<\Psi<2h'\}\cap\Cell}\mathcal J_\omega$ by Lemma~\ref{lem:coarea-level-cutoff} with
$\varphi=\mathbf 1_{(h,h')}$. Finally $\tilde G=G_\omega$ almost everywhere.
\end{proof}

\begin{lemma}[The level terms at the center]
\label{lem:ee-center-limit}
\uses{def:ee-level-term}
\lean{EulerDipoleEigenvalues.ee_center_limit}
\leanok
For $\omega\in L^2$ and every null set $N$ there are $h_j\in(0,1)\setminus N$ with $h_j\to1$
and $G_\omega(h_j)\to0$.
\end{lemma}

\begin{proof}
\uses{lem:orbit-coordinates-area, lem:continuous-representative, lem:lift-representatives,
  lem:level-integral-congr, lem:period-integral, lem:period-orbit}
\leanok
Let $F=|\psi|^2+|\psi_x|^2+|\psi_y|^2+|\omega|^2$, integrable on the cell, and
$g(h)=\oint_{\{\Psi=2h\}}F\,dt$. By Lemma~\ref{lem:orbit-coordinates-area},
$2\int_0^1g\,dh=\int_\Cell F<\infty$. Hence for every null set $N$ there are $h_j\to1$
outside $N$ with $(1-h_j)g(h_j)\to0$: otherwise $g\ge c/(1-h)$ almost everywhere near $1$,
which is not integrable.

It remains to bound $G_\omega(h)$ for almost every $h\in(1/2,1)$. Since $\psi$ agrees almost
everywhere with its continuous representative (Lemma~\ref{lem:lift-representatives}),
$|\psi|\le M=\sup|\tilde\psi|$ almost everywhere on almost every level
(Lemma~\ref{lem:level-integral-congr}). Pointwise $|\psi_n|^2,|\psi_t|^2\le2F$ and
$|\operatorname{Re}(\omega\bar\psi)|\le F$, so the second term of $G_\omega$ is at most
$10(1-h)g(h)$. For the first, $\rho^2=(1-\cos x)(1+\cos x)+(1-\cos y)(1+\cos y)
\le2(2-\Psi)=4(1-h)$ on the level, exactly and not only near the center, and
$S(h)=8L(h)\le8/h\le16$ (Lemma~\ref{lem:period-integral}). By Cauchy--Schwarz
\[
  \Bigl|\oint\operatorname{Re}(\psi\overline{\psi_n})\,\rho\,dt\Bigr|
  \le M\Bigl(\oint\rho^2dt\Bigr)^{1/2}\Bigl(\oint|\psi_n|^2dt\Bigr)^{1/2}
  \le M\bigl(128(1-h)g(h)\bigr)^{1/2}.
\]
Both bounds tend to $0$ along $h_j$.
\end{proof}

\begin{lemma}[Truncated virial identity]
\label{lem:ee-truncated-virial}
\uses{def:ee-bulk-density, def:eigen-problem, def:level-integral, def:normal-tangent, def:stream-lifts}
\lean{EulerDipoleEigenvalues.ee_truncated_virial}
\leanok
Let $\lambda\ne0$. For almost every $h\in(0,1)$, put $\epsilon=2h$. The density $\mathcal J_\omega$ is
integrable on $\{p\in\Cell:\Psi(p)>\epsilon\}$. With $n$ pointing out of this set, the level
integrals in flow time $dt$ over $\{\Psi=\epsilon\}$ and $d\ell=\rho\,dt$ satisfy
\[
  \int_{\{\Psi>\epsilon\}}\mathcal J
  =\oint_{\{\Psi=\epsilon\}}\operatorname{Re}(\psi\overline{\psi_n})\,d\ell
   +(2-\epsilon)\left(\oint|\psi_n|^2\,dt-\oint|\psi_t|^2\,dt
     -\oint\operatorname{Re}(\omega\bar\psi)\,dt\right)
\]
(manuscript \texttt{(ee:truncated-virial)}).
\end{lemma}

\begin{proof}
\uses{lem:ee-bulk-integrable, lem:ee-annular-virial, lem:ee-center-limit, lem:cos-diagonal-null}
\leanok
Let $h\in(0,1)$ lie outside the null set $N$ of Lemma~\ref{lem:ee-annular-virial}.
Lemma~\ref{lem:ee-center-limit} gives $h'_j\to1$ outside $N$ with $G_\omega(h'_j)\to0$, so
$\int_{\{2h<\Psi<2h'_j\}}\mathcal J=G_\omega(h)-G_\omega(h'_j)\to G_\omega(h)$. On the other hand
these integrals tend to $\int_{\{\Psi>2h\}}\mathcal J$ by dominated convergence
(Lemma~\ref{lem:ee-bulk-integrable}), since $\{\Psi=2\}\subset\{\cos x=\cos y\}$ is null
(Lemma~\ref{lem:cos-diagonal-null}).
\end{proof}

\begin{remark}[Digest note: integrability in the truncated identity]
An equality of total Bochner integrals alone does not justify positivity:
Lean assigns integral zero to nonintegrable functions. The truncated
identity therefore includes integrability and $\lambda\ne0$. Only
integrability on each truncated cell is needed, not at the saddles;
Lemma~\ref{lem:nonnegative-exhaustion} gives the final passage to the cell.
\end{remark}

\section{Even--even gradient traces}
\label{sec:ee-gradient-limits}

\begin{remark}[Parity and traces]
Almost-everywhere parity of the raw derivatives is not a pointwise
boundary condition. Lemma~\ref{lem:arc-trace-identification} uses
continuous line representatives to evaluate this parity on the fixed
separatrix arc and obtain the zero normal trace.
\end{remark}

\begin{lemma}[Energy of the arc traces]
\label{lem:ee-arc-energy}
\uses{def:sectors, def:stream-lifts, def:cell-arcs}
\lean{EulerDipoleEigenvalues.ee_arc_energy, EulerDipoleEigenvalues.lintegral_arcTrace_lt_top}
\leanok
Let $\psi=\Delta^{-1}\omega$, $\omega\in X$, be even under $S_1,S_2$, and let
$g_k=\tilde\psi\circ\gamma_k$ be its traces on the four arcs. Then
$\int\cosh^2s\,|g_k'|^2\,ds<\infty$. Equivalently (for $k=0$), with $f(a)=\tilde\psi(a,\pi/2)$,
$\int_{-\pi/2}^{\pi/2}|f'(a)|^2\,da/\cos a<\infty$.
\end{lemma}

\begin{proof}
\uses{lem:arc-trace-identification, lem:line-representatives, lem:trace-gagliardo, lem:lift-sobolev-local,
  lem:continuous-stream-symmetries, lem:local-odd-hardy}
\leanok
The arc $\gamma_k$ for $\omega$ is the arc $\gamma_0$ for $\omega\circ R^k$, again even under $S_1,S_2$
(Lemma~\ref{lem:continuous-stream-symmetries}), so take $k=0$. By
Lemma~\ref{lem:arc-trace-identification}(4) the energy is $\int|f'|^2da/\cos a$. By
Lemma~\ref{lem:trace-gagliardo} applied to $\psi_a=\psi_x+\psi_y\in H^1$
(Lemma~\ref{lem:lift-sobolev-local}), its line representative has a trace on $b=\pi/2$ with finite
Gagliardo seminorm on $(-\pi,\pi)$, and this trace is $f'$ almost everywhere
(Lemmas~\ref{lem:line-representatives}(3) and~\ref{lem:arc-trace-identification}(2)). Since
$\tilde\psi$ is even under $S_1$ and periodic, $f(\pi-a)=f(a)=f(-\pi-a)$ for every $a$, so
$f'(\pi/2+\tau)=-f'(\pi/2-\tau)$ and $f'(-\pi/2+\tau)=-f'(-\pi/2-\tau)$ for every $\tau$.
Lemma~\ref{lem:local-odd-hardy} with $\delta=\pi/2$ bounds $\int_0^{\pi/2}|f'(\pm(\pi/2-u))|^2du/u$ by
twice the Gagliardo seminorm, and $\cos a=\sin(\pi/2-|a|)\ge\frac2\pi(\pi/2-|a|)$ (Jordan).
\end{proof}

\begin{lemma}[Level traces at the separatrix]
\label{lem:ee-level-traces}
\uses{def:sectors, def:normal-tangent, def:stream-lifts, def:level-integral, def:cell-arcs}
\lean{EulerDipoleEigenvalues.ee_level_traces, EulerDipoleEigenvalues.ee_quarter_traces,
  EulerDipoleEigenvalues.tendsto_quarter_integral,
  EulerDipoleEigenvalues.continuousAt_line_normalDeriv}
\leanok
For $\psi$ as in Lemma~\ref{lem:ee-arc-energy}, as $h\to0$ outside a null set,
\[
  \oint|\psi_n|^2\,dt\to0,\qquad
  \oint|\psi_t|^2\,dt\to\sum_{k=0}^{3}\int_\R\tfrac12\cosh^2s\,|g_k'(s)|^2\,ds
\]
over the levels $\{\Psi=2h\}$ (manuscript \texttt{(ee:normal-trace-limit)},
\texttt{(ee:tangent-trace-limit)}).
\end{lemma}

\begin{proof}
\uses{lem:ee-quarter-bound, lem:arc-trace-identification, lem:line-representatives,
  lem:quarter-orbit-decomposition, lem:quarter-orbit-diagonal, lem:quarter-null-sets,
  lem:arc-diagonal, lem:normal-tangent-rotation, lem:continuous-stream-symmetries}
\leanok
By Lemma~\ref{lem:quarter-orbit-decomposition} the level integrals are sums of four quarter
integrals, and the $k$-th quarter of $\psi$ is the $0$-th quarter of the stream function of
$\omega\circ R^k$, again even under $S_1,S_2$ (Lemmas~\ref{lem:normal-tangent-rotation}
and~\ref{lem:continuous-stream-symmetries}; the almost-everywhere identities hold on almost
every level by Lemma~\ref{lem:quarter-null-sets}). So consider $k=0$. Let $G_x,G_y$ be line
representatives of $\psi_x,\psi_y$; by Lemma~\ref{lem:quarter-null-sets} they agree with the raw
$\psi_x,\psi_y$ at $(a,\beta_h(a))$ for almost every $a$, for $h$ outside a null set. By
Lemma~\ref{lem:quarter-orbit-diagonal} the quarter integral of $F=|\psi_n|^2$ or $|\psi_t|^2$ is
\[
  \int_{-\pi/2}^{\pi/2}\mathbf 1_{|a|<a_h}\,F(a,\beta_h(a))\,\frac{da}{\sqrt{\cos^2a-h^2}} .
\]
For $h\le\frac14$ the integrand is bounded by the integrable $\Phi$ of
Lemma~\ref{lem:ee-quarter-bound}, uniformly in $h$. For almost every $a$, as $h\to0$:
$a_h\to\pi/2>|a|$, $\beta_h(a)=\arccos(h/\cos a)\to\pi/2$, $G_x,G_y$ are continuous in $b$ on the
line through $a$ (Lemma~\ref{lem:line-representatives}), $\rho\ne0$ at $(a,\pi/2)$, and the weight
tends to $1/\cos a$. By dominated convergence the quarter integral tends to
$\int_{-\pi/2}^{\pi/2}F(a,\pi/2)\,da/\cos a$. For $F=|\psi_n|^2$ this is $0$, and for
$F=|\psi_t|^2$ it is $\int\frac12\cosh^2s\,|g'|^2ds$ (Lemmas~\ref{lem:arc-diagonal}
and~\ref{lem:arc-trace-identification}).
\end{proof}

\begin{remark}[Digest note: level traces]
The bound of Lemma~\ref{lem:ee-quarter-bound} holds on the entire quarter
for $h\le\frac14$, so one dominated-convergence argument replaces a
separate middle/corner split. Fubini in $(a,h)$
(Lemma~\ref{lem:quarter-null-sets}) handles changes of representatives;
these quadratic trace limits need no coarea formula.
\end{remark}

\begin{lemma}[Traces on the arc]
\label{lem:arc-trace-identification}
\uses{def:sectors, def:stream-lifts, def:normal-tangent, def:separatrix-arc,
  def:line-representative}
\lean{EulerDipoleEigenvalues.ee_arc_trace_identification,
  EulerDipoleEigenvalues.lintegral_arc_energy_eq}
\leanok
Let $\psi$ be as in Lemma~\ref{lem:ee-arc-energy}, and let $G_x,G_y$ be line representatives
(Definition~\ref{def:line-representative}) of $\psi_x,\psi_y$ on the diagonal rectangle
$\{|a|<\pi,\ 0<b<\pi\}$. Then:
\begin{enumerate}
  \item for almost every $a$, $G_x(a,\pi/2)=G_y(a,\pi/2)$, that is the trace of
        $\psi_b=\psi_x-\psi_y$ vanishes, and hence so does the normal derivative
        $(G_x\sin x+G_y\sin y)/\rho$ at $(a,\pi/2)$;
  \item for almost every $a$, $a\mapsto\tilde\psi(a,\pi/2)$ has derivative
        $\psi_a(a,\pi/2):=G_x(a,\pi/2)+G_y(a,\pi/2)$;
  \item for almost every $s$, $g=\tilde\psi\circ\gamma$ has derivative
        $g'(s)=T(\gamma(s))$ with $T=\sin y\,G_x-\sin x\,G_y$, and
        $|\psi_t|^2=\frac12\cosh^2s\,|g'(s)|^2$ at $\gamma(s)$, where $\psi_t=T/\rho$.
  \item With $f(a)=\tilde\psi(a,\pi/2)$,
        $\int\cosh^2s\,|g'(s)|^2\,ds=\int_{-\pi/2}^{\pi/2}|f'(a)|^2\,da/\cos a$ (as Lebesgue
        integrals in $[0,\infty]$).
\end{enumerate}
\end{lemma}

\begin{proof}
\uses{lem:stream-derivative-parities, lem:line-representatives, lem:trace-h2,
  lem:lift-sobolev-local, lem:lift-representatives, lem:diagonal-calculus, lem:separatrix-arc,
  lem:arc-diagonal, lem:normal-tangent-frame}
\leanok
(1) Combined with a period, $S_2$ is the reflection $\sigma(x,y)=(y+\pi,x-\pi)$, which is
$b\mapsto\pi-b$ at fixed $a$. Since $\psi$ is even under $S_2$, $\psi_x\circ\sigma=\psi_y$ and
$\psi_y\circ\sigma=\psi_x$ almost everywhere (Lemma~\ref{lem:stream-derivative-parities}), so
$D=G_x-G_y$ satisfies $D(a,\pi-b)=-D(a,b)$ for almost every $b$ on almost every line
(Lemma~\ref{lem:diagonal-calculus}). Both sides are continuous in $b$ on almost every line
(Lemma~\ref{lem:line-representatives}), so the identity holds for all $b$, and at $b=\pi/2$ it
gives $D=0$. On $b=\pi/2$, $\sin x=\cos a$ and $\sin y=-\cos a$, so the normal numerator is
$\cos a\,D=0$.

(2) Lemma~\ref{lem:trace-h2}, applied to the continuous $H^2$ function $\tilde\psi$
(Lemmas~\ref{lem:lift-sobolev-local} and~\ref{lem:lift-representatives}), gives a representative
$H$ of $\psi_x+\psi_y$, continuous on almost every line, such that $a\mapsto\tilde\psi(a,\pi/2)$
has derivative $H(a,\pi/2)$ almost everywhere. By
Lemma~\ref{lem:line-representatives}(3), $H=G_x+G_y$ on almost every whole line, in particular
at $b=\pi/2$.

(3) $g=f\circ\phi$ with $f(a)=\tilde\psi(a,\pi/2)$ and $\phi(s)=\alpha(s)-\pi/2$,
$\phi'(s)=-\sech s$ (Lemma~\ref{lem:separatrix-arc}). The inverse
$a\mapsto-\log\tan((a+\pi/2)/2)$ of $\phi$ is differentiable, so it maps null sets to null
sets, and for almost every $s$ the point $\phi(s)$ is one where (2) holds. There the chain rule
gives $g'(s)=-\sech s\,(G_x+G_y)(\gamma(s))$, which is $T(\gamma(s))$ since on the arc
$\sin x=\sech s$ and $\sin y=-\sech s$. Finally $\rho(\gamma(s))^2=2\sech^2s$.

(4) By (3), $\cosh^2s\,|g'|^2=2|\psi_t|^2$ at $\gamma(s)$; by Lemma~\ref{lem:arc-diagonal} this
integrates to $\int 2|\psi_t(a,\pi/2)|^2da/\cos a$, and on $b=\pi/2$,
$2|\psi_t|^2=2(|G_x|^2+|G_y|^2)=|G_x+G_y|^2=|f'|^2$ by (1), (2) and
Lemma~\ref{lem:normal-tangent-frame}.
\end{proof}

\begin{remark}[Choice of line representatives]
The arc traces are independent of the chosen line representatives by
Lemma~\ref{lem:line-representatives}(3). Their continuity in $b$ suffices
to evaluate reflection parity on the arc; Lemma~\ref{lem:trace-h2}
identifies the tangential trace with the derivative of the stream trace.
\end{remark}

\begin{lemma}[A dominating function on the quarter orbits]
\label{lem:ee-quarter-bound}
\uses{def:sectors, def:stream-lifts, def:normal-tangent, def:line-representative}
\lean{EulerDipoleEigenvalues.ee_quarter_bound}
\leanok
Let $\psi$ and $G_x,G_y$ be as in Lemma~\ref{lem:arc-trace-identification}. There is $\Phi\ge0$ with
$\int_{-\pi/2}^{\pi/2}\Phi<\infty$ such that for $h\in(0,1)$ and almost every $|a|<a_h$ with
$|\sin a|\ge\frac12$ or $h\le\frac14$,
\[
  \bigl(|\psi_n|^2+|\psi_t|^2\bigr)(a,\beta_h(a))\,\frac1{\sqrt{\cos^2a-h^2}}\le\Phi(a),
\]
with $\psi_n,\psi_t$ formed from $G_x,G_y$.
\end{lemma}

\begin{proof}
\uses{lem:ee-arc-energy, lem:arc-trace-identification, lem:line-representatives,
  lem:normal-tangent-frame, lem:segment-estimate, lem:quarter-height-limit,
  lem:quarter-weight}
\leanok
The weight is at most $2/\cos a$ (Lemma~\ref{lem:quarter-weight}). By
Lemma~\ref{lem:normal-tangent-frame}, $|\psi_n|^2+|\psi_t|^2=|G_x|^2+|G_y|^2$ on the quarter orbit
($\rho\ne0$ there since $\cos a>0$ and $\sin\beta_h(a)>0$). Write
$|G(a,\beta_h(a))|^2\le2|G(a,\beta_h(a))-G(a,\pi/2)|^2+2|G(a,\pi/2)|^2$ for $G=G_x,G_y$.

Differences: by the segment estimate (Lemma~\ref{lem:segment-estimate}) with $w=1/\cos a$,
$J=(0,\pi)$, and $\pi/2-\beta_h(a)\le\frac\pi2h/\cos a$ (Lemma~\ref{lem:quarter-height-limit}),
$w\,|\beta_h-\pi/2|\le\frac\pi2h/\cos^2a\le\frac\pi2$, so the difference terms are bounded by
$\frac\pi2\int_0^\pi(|\partial_bG_x|^2+|\partial_bG_y|^2)\,db$, which is integrable in $a$
(Lemma~\ref{lem:line-representatives}).

Traces: by Lemma~\ref{lem:arc-trace-identification},
$(|G_x|^2+|G_y|^2)(a,\pi/2)=\frac12|f'(a)|^2$ with $f(a)=\tilde\psi(a,\pi/2)$, and
$\int_{-\pi/2}^{\pi/2}|f'|^2da/\cos a<\infty$ (Lemma~\ref{lem:ee-arc-energy}). So
$\Phi=4\bigl(\frac\pi2(K_x+K_y)+|f'|^2/(2\cos a)\bigr)$ works, with $K$ the line integrals of
$|\partial_bG|^2$. The pointwise form of the segment estimate is used since $G(a,\beta_h(a))$ need
not be measurable in $a$.
\end{proof}

\begin{lemma}[The first-order trace term]
\label{lem:ee-first-trace-limit}
\uses{def:normal-tangent, def:stream-lifts, def:level-integral}
\lean{EulerDipoleEigenvalues.ee_first_trace_limit,
  EulerDipoleEigenvalues.abs_levelIntegral_first_le,
  EulerDipoleEigenvalues.abs_levelIntegral_first_le_of_bounded,
  EulerDipoleEigenvalues.levelIntegral_gradNorm_sq,
  EulerDipoleEigenvalues.levelIntegral_gradNorm_sq_le}
\leanok
If $\oint|\psi_n|^2\,dt\to0$ as $h\to0$ outside a null set $N$, then
$\oint\operatorname{Re}(\psi\overline{\psi_n})\,d\ell\to0$ outside a null set containing $N$
(manuscript \texttt{(ee:first-trace-limit)}).
\end{lemma}

\begin{proof}
\uses{lem:continuous-representative, lem:lift-representatives, lem:level-integral-congr,
  lem:orbit-coordinate-gradients, lem:orbit-time-integrals, lem:orbit-coordinates-area,
  lem:period-orbit}
\leanok
On almost every level $|\psi|\le M=\sup|\tilde\psi|$, and $d\ell=\rho\,dt$, so by
Cauchy--Schwarz $|\oint\operatorname{Re}(\psi\overline{\psi_n})\,d\ell|\le
M(\oint\rho^2dt)^{1/2}(\oint|\psi_n|^2dt)^{1/2}$.
The same bound holds for arbitrary measurable $f,z$ on the level with
$|f|\le M$ almost everywhere and $z\in L^2(dt)$, replacing
$\psi,\psi_n$ by $f,z$; no derivative relation is required.
Although the period $S(h)\to\infty$,
$\oint\rho^2\,dt$ is bounded: $\rho\le Cd$ with $d$ the distance to the saddles, and
$\int d^2\,dt\le C$ on each quarter orbit (Lemma~\ref{lem:orbit-time-integrals}); $\rho^2$ is
invariant under the quarter rotation, so $\oint\rho^2dt$ is four quarter integrals.
(Equivalently, $\oint\rho^2dt=\oint\nabla\Psi\cdot n\,d\ell=\int_{\{\Psi>2h\}}\Psi$ by the divergence
theorem.)
\end{proof}

\begin{lemma}[Weighted gradient traces]
\label{lem:ee-gradient-traces}
\uses{lem:ee-arc-energy, lem:ee-level-traces}
\lean{EulerDipoleEigenvalues.ee_gradient_traces}
\leanok
For $\psi$ as in Lemma~\ref{lem:ee-arc-energy}, $\int\cosh^2s\,|g_k'|^2\,ds<\infty$, and as
$h\to0$ outside a null set,
\[
  \oint|\psi_n|^2\,dt\to0,\qquad
  \oint|\psi_t|^2\,dt\to\sum_{k=0}^{3}\int_\R\tfrac12\cosh^2s\,|g_k'(s)|^2\,ds,\qquad
  \oint\operatorname{Re}(\psi\overline{\psi_n})\,d\ell\to0
\]
over the levels $\{\Psi=2h\}$ (manuscript \texttt{(ee:gradient-traces)}).
\end{lemma}

\begin{proof}
\uses{lem:ee-arc-energy, lem:ee-level-traces, lem:ee-first-trace-limit}
\leanok
Lemmas~\ref{lem:ee-arc-energy}, \ref{lem:ee-level-traces} and~\ref{lem:ee-first-trace-limit},
with the union of the two null sets.
\end{proof}

\begin{remark}[Mechanism of the gradient limits]
Reflection parity gives the zero normal limit; the $H^2$ trace theorem
identifies the tangential limit with $\frac12\cosh^2s\,|g'|^2$.
The local odd Hardy inequality supplies the integrable arc energy.
The first-order term is separate: it needs $\oint\rho^2\,dt$ bounded,
not the period bounded, since the period diverges at the separatrix.
All level limits are taken outside their combined exceptional null set.
\end{remark}

\section{One-dimensional boundary positivity}
\label{sec:ee-boundary-positivity}

The forms, weighted domain and algebra below do not depend on reflection
parity. The eigenfunction application begins with
Lemma~\ref{lem:arc-trace-ac}.

\subsection{Forms, domain and classical identities}

\begin{definition}[Boundary form]
\label{def:boundary-form}
\uses{def:boundary-form-xi}
\lean{EulerDipoleEigenvalues.flowBoundaryForm,
  EulerDipoleEigenvalues.flowValue,
  EulerDipoleEigenvalues.flowFirst,
  EulerDipoleEigenvalues.flowSecond}
\leanok
For a function $z$ on $\R$ with $P=-i\partial_s$, let
(manuscript \texttt{(ee:trace-form)})
\[
  \mathcal F_\lambda(z)=\frac12\int_\R\cosh^2s\,|\partial_s(P-\lambda)z|^2\,ds
  -\operatorname{Re}\int_\R Pz\,\overline{(P-\lambda)z}\,ds .
\]
In the variable $\xi=\tanh s$, with $w=1-\xi^2$ and $F(\xi)=z(s)$, we also write
$\mathcal F_\lambda(F)$, and $G=-iwF'-\lambda F$ corresponds to $(P-\lambda)z$.
Given $F,D,E$ on $(-1,1)$, the corresponding functions of $s$ are
$z=F\circ\tanh$, $z_1=(wD)\circ\tanh$ and $z_2=(w(wE-2\xi D))\circ\tanh$.
As for $\mathcal B_\lambda$, Lean takes the derivatives $z_1=z'$, $z_2=z_1'$
as explicit arguments, and integrates the real part in the second term.
\end{definition}

\begin{definition}[Finite-interval boundary form]
\label{def:boundary-form-xi}
\lean{EulerDipoleEigenvalues.boundaryWeight, EulerDipoleEigenvalues.boundaryForm}
\leanok
For $\lambda\in\R$ and $F,D,E:\R\to\C$, put $w(\xi)=1-\xi^2$ and
\[
  \mathcal B_\lambda(F,D,E)
  =\frac12\int_{-1}^1|wE-2\xi D-i\lambda D|^2\,d\xi
   -\int_{-1}^1w|D|^2\,d\xi
   +\lambda\int_{-1}^1\operatorname{Im}(D\bar F)\,d\xi.
\]
The derivatives are explicit arguments; their relation to $F$ is imposed in
subsequent statements. Lean uses total Bochner integrals in this definition;
applications below ensure that the integrands are integrable.
\end{definition}

\begin{definition}[Classical boundary jet]
\label{def:boundary-jet}
\lean{EulerDipoleEigenvalues.BoundaryJet}
\leanok
A classical boundary jet is a triple $(F,D,E)$ of functions $\R\to\C$,
continuous on $[-1,1]$, with $F'(\xi)=D(\xi)$ and $D'(\xi)=E(\xi)$ at every
$\xi\in[-1,1]$. The derivatives at the endpoints are ordinary two-sided
derivatives of the functions on $\R$.
\end{definition}

\begin{definition}[Weighted second-order boundary domain]
\label{def:weighted-boundary-domain}
\uses{def:h1-interval, def:boundary-form-xi}
\lean{EulerDipoleEigenvalues.HasWeightedBoundaryH2,
  EulerDipoleEigenvalues.MemWeightedBoundaryH2}
\leanok
Write $I=(-1,1)$ and $w(\xi)=1-\xi^2$. A weighted boundary triple $(F,D,E)$
consists of $H^1(I)$ data $(F,D)$, with $D$ absolutely continuous on every
compact subinterval of $I$, $D'=E$ almost everywhere on $I$, and
$wE\in L^2(I)$. Let $\mathcal D_w$ be the set of $F$ admitting such $D,E$.
There is no global $L^2$ condition on $E$, no endpoint trace condition on
$D$, and no restriction on raw endpoint or outside values.
\end{definition}

\begin{lemma}[Classical inclusion and traces]
\label{lem:boundary-jet-weighted}
\uses{def:boundary-jet, def:h1-interval, def:weighted-boundary-domain, lem:h1-interval-endpoints}
\lean{EulerDipoleEigenvalues.BoundaryJet.hasWeightedBoundaryH2,
  EulerDipoleEigenvalues.BoundaryJet.weighted_extension_eq}
\leanok
Every classical boundary jet $(F,D,E)$ is a weighted boundary triple, and
its canonical $H^1$ extension agrees with $F$ on $[-1,1]$. In particular
the trace values are its ordinary endpoint values.
\end{lemma}

\begin{proof}
\uses{lem:h1-interval-endpoints}
\leanok
Continuous functions on a compact interval are bounded and square
integrable. The continuous derivatives $D$ and $E$ give Lipschitz, hence
absolutely continuous, restrictions of $F$ and $D$ to compact subintervals.
The weighted second derivative is continuous on $[-1,1]$. Uniqueness of the
continuous extension identifies the traces.
\end{proof}

\begin{lemma}[Linear and affine stability]
\label{lem:weighted-boundary-linear}
\uses{def:h1-interval, def:weighted-boundary-domain}
\lean{EulerDipoleEigenvalues.HasWeightedBoundaryH2.memH1,
  EulerDipoleEigenvalues.HasWeightedBoundaryH2.add,
  EulerDipoleEigenvalues.HasWeightedBoundaryH2.smul,
  EulerDipoleEigenvalues.HasWeightedBoundaryH2.sub,
  EulerDipoleEigenvalues.hasWeightedBoundaryH2_affine,
  EulerDipoleEigenvalues.HasWeightedBoundaryH2.subAffine}
\leanok
Weighted boundary triples are closed under addition, subtraction and complex
scalar multiplication, with the corresponding operations on all three
components. Every affine function $a+b\xi$ belongs to $\mathcal D_w$ with
derivatives $b,0$. Consequently, subtracting it from a triple $(F,D,E)$
gives the triple $(F-a-b\xi,D-b,E)$. Also $\mathcal D_w\subset H^1(I)$.
\end{lemma}

\begin{proof}
\uses{lem:boundary-jet-weighted}
\leanok
Absolute continuity, a.e.\ differentiation and $L^2$ membership are linear.
Affine functions are classical boundary jets.
\end{proof}

\begin{lemma}[Interior regularity of the first derivative]
\label{lem:weighted-boundary-local}
\uses{def:h1-interval, def:weighted-boundary-domain}
\lean{EulerDipoleEigenvalues.HasWeightedBoundaryH2.memLp_second_local,
  EulerDipoleEigenvalues.HasWeightedBoundaryH2.first_hasIntervalH1}
\leanok
For every weighted boundary triple $(F,D,E)$ and $-1<c<d<1$, one has
$E\in L^2(c,d)$ and $D\in H^1(c,d)$ with a.e.\ derivative $E$.
\end{lemma}

\begin{proof}
\leanok
The inverse weight $1/w$ is continuous and bounded on $[c,d]$, so
$E=(1/w)(wE)$ is square integrable there. Combine this with the local
absolute continuity of $D$ and the given a.e.\ derivative relation.
\end{proof}

\begin{remark}[Representative model and endpoint singularities]
The domain uses locally absolutely continuous representatives for the
first derivative, not global $H^1$ regularity or endpoint traces of it.
For example, $F(\xi)=(1-\xi)^{3/2}$ has
$E=\tfrac34(1-\xi)^{-1/2}$ and
$wE=\tfrac34(1+\xi)\sqrt{1-\xi}$. Thus $F\in\mathcal D_w$ although
$\int|E|^2$ diverges logarithmically: the domain is larger than $H^2$.
\end{remark}

\begin{lemma}[Boundary form in the $\xi$ variable]
\label{lem:boundary-form-xi}
\uses{def:boundary-form, def:boundary-form-xi, def:weighted-boundary-domain}
\lean{EulerDipoleEigenvalues.HasWeightedBoundaryH2.flowBoundaryForm_eq,
  EulerDipoleEigenvalues.hasDerivAt_flowValue,
  EulerDipoleEigenvalues.hasDerivAt_flowFirst,
  EulerDipoleEigenvalues.integral_tanh_subst,
  EulerDipoleEigenvalues.boundaryWeight_tanh}
\leanok
Let $(F,D,E)$ be a weighted boundary triple. Then
\[
  \mathcal F_\lambda(F)=\frac12\int_{-1}^1|wF''-2\xi F'-i\lambda F'|^2\,d\xi
  -\int_{-1}^1w|F'|^2\,d\xi+\lambda\operatorname{Im}\int_{-1}^1F'\bar F\,d\xi ,
\]
that is, $\mathcal F_\lambda(F)=\mathcal B_\lambda(F,D,E)$, where
$\mathcal F_\lambda(F)$ is computed from $z=F\circ\tanh$ and the derivatives
$z_1=(wD)\circ\tanh$, $z_2=(w(wE-2\xi D))\circ\tanh$. These are the derivatives
of $z$ and $z_1$ at every $s$ where $F'=D$, respectively $D'=E$, holds at
$\tanh s$ (manuscript \texttt{(ee:trace-xi)}).
\end{lemma}

\begin{proof}
\uses{lem:separatrix-arc, lem:weighted-boundary-integrable}
\leanok
$\partial_s=w\partial_\xi$, $\cosh^2s=1/w$, $ds=d\xi/w$, and
$G'=-i(wF''-2\xi F'-i\lambda F')$. Pointwise, the first integrand is
$\cosh^{-2}s\,|wE-2\xi D-i\lambda D|^2$ and the second is
$\cosh^{-2}s\,(w|D|^2-\lambda\operatorname{Im}(D\bar F))$ at $\xi=\tanh s$. The
substitution $\xi=\tanh s$ is a change of variables for an injective map
onto $(-1,1)$ and needs no integrability; integrability
(Lemma~\ref{lem:weighted-boundary-integrable}) is used only to split the
second integral.
\end{proof}

\begin{lemma}[Almost-everywhere flow derivatives]
\label{lem:weighted-flow-derivatives}
\uses{def:boundary-form, def:h1-interval, def:weighted-boundary-domain}
\lean{EulerDipoleEigenvalues.ae_comp_tanh,
  EulerDipoleEigenvalues.HasWeightedBoundaryH2.ae_hasDerivAt_flowValue,
  EulerDipoleEigenvalues.HasWeightedBoundaryH2.ae_hasDerivAt_flowFirst,
  EulerDipoleEigenvalues.HasWeightedBoundaryH2.continuous_flowFirst}
\leanok
Composition with $\tanh$ preserves almost-everywhere statements on $\R$.
For a weighted boundary triple $(F,D,E)$, put
$z_0=F\circ\tanh$, $z_1=(wD)\circ\tanh$ and
$z_2=(w(wE-2\xi D))\circ\tanh$. Then $z_0'=z_1$ and $z_1'=z_2$ almost everywhere,
and $z_1$ is continuous.
\end{lemma}

\begin{proof}
\uses{lem:boundary-form-xi}
\leanok
The preimage of a null set under $\tanh$ is its intersection with $(-1,1)$ mapped by
$\operatorname{artanh}$. The latter is differentiable there and sends null sets to null
sets. Pull back the two almost-everywhere derivative relations of the weighted domain
and apply the pointwise chain rules of Lemma~\ref{lem:boundary-form-xi}.
Local absolute continuity makes $D$ continuous in the interior, hence $z_1$ continuous.
\end{proof}

\begin{lemma}[Classical completion of squares]
\label{lem:trace-completion-classical}
\uses{def:boundary-form-xi, def:boundary-jet}
\lean{EulerDipoleEigenvalues.BoundaryJet.completion,
  EulerDipoleEigenvalues.BoundaryJet.sum_of_squares}
\leanok
For a classical boundary jet $(F,D,E)$ and real $\lambda$, put
$A=wE-\xi D+F-i\lambda D$. Then
\[
  \mathcal B_\lambda(F,D,E)
  =\frac12\int_{-1}^1|A|^2\,d\xi
   +\frac12\int_{-1}^1|D|^2\,d\xi
   -\frac12\bigl(|F(1)|^2+|F(-1)|^2\bigr).
\]
In particular, the last term vanishes if $F(\pm1)=0$.
\end{lemma}

\begin{proof}
\leanok
The difference between the original density and
$\tfrac12|A|^2+\tfrac12|D|^2$ is $K'$, where
\[
  K=-\tfrac12 w\xi|D|^2-w\operatorname{Re}(D\bar F)-\tfrac12\xi|F|^2.
\]
This is a pointwise differentiation and polynomial identity, including all
terms involving $\lambda$. Integrate by the fundamental theorem of calculus.
Since $w(\pm1)=0$, $K(1)-K(-1)$ is exactly the displayed endpoint correction.
\end{proof}

\begin{lemma}[Sum of squares, smooth case]
\label{lem:trace-sos-smooth}
\uses{lem:boundary-form-xi, def:boundary-jet}
\lean{EulerDipoleEigenvalues.BoundaryJet.sum_of_squares}
\leanok
For a classical boundary jet $(F,D,E)$ with $F(-1)=F(1)=0$ and real $\lambda$,
\[
  \mathcal F_\lambda(F)=\frac12\int_{-1}^1|wE-\xi D+F-i\lambda D|^2\,d\xi
  +\frac12\int_{-1}^1|D|^2\,d\xi.
\]
In particular this applies to $F\in C_c^\infty(-1,1)$ with $D=F'$ and $E=F''$
(manuscript \texttt{(ee:trace-sos)}).
\end{lemma}

\begin{proof}
\uses{lem:trace-completion-classical}
\leanok
Apply the classical completion to $(F,D,E)$. The endpoint correction vanishes.
For a compactly supported smooth $F$, extend by zero to $\R$ and take
$(D,E)=(F',F'')$.
\end{proof}

\begin{remark}
The classical nodes above are the special case of the weighted nodes below
in which the endpoint flux can be evaluated directly at $\pm1$.
\end{remark}

\subsection{The weighted domain}

The classical argument extends to $\mathcal D_w$ without any approximation:
integrate the flux identity on interior intervals, where a weighted triple
is an ordinary first-order pair, and let the interval exhaust $(-1,1)$. The
only point is that the weighted part of the flux vanishes along suitable
sequences, which follows from $D\in L^2$ alone.

\begin{lemma}[Pointwise flux identity]
\label{lem:boundary-flux}
\uses{def:boundary-form-xi}
\lean{EulerDipoleEigenvalues.boundaryFlux,
  EulerDipoleEigenvalues.boundaryFluxDerivative,
  EulerDipoleEigenvalues.boundary_density_completion,
  EulerDipoleEigenvalues.hasDerivAt_boundaryFlux}
\leanok
Let $\lambda\in\R$, $\xi\in\R$, and let $F,D,E:\R\to\C$ with
$F'(\xi)=D(\xi)$ and $D'(\xi)=E(\xi)$. Put
\[
  K=-\tfrac12 w\xi|D|^2-w\operatorname{Re}(D\bar F)-\tfrac12\xi|F|^2,\qquad
  A=wE-\xi D+F-i\lambda D.
\]
Then $K$ is differentiable at $\xi$ and, at $\xi$,
\[
  \tfrac12|wE-2\xi D-i\lambda D|^2-w|D|^2+\lambda\operatorname{Im}(D\bar F)
  =\tfrac12|A|^2+\tfrac12|D|^2+K'.
\]
\end{lemma}

\begin{proof}
\leanok
Product rule for $K$, then a polynomial identity in the real and imaginary
parts of $F(\xi),D(\xi),E(\xi)$.
\end{proof}

\begin{lemma}[Integrability on the weighted domain]
\label{lem:weighted-boundary-integrable}
\uses{def:h1-interval, def:weighted-boundary-domain, def:boundary-form-xi}
\lean{EulerDipoleEigenvalues.HasWeightedBoundaryH2.memLp_boundarySquare,
  EulerDipoleEigenvalues.HasWeightedBoundaryH2.memLp_formSquare,
  EulerDipoleEigenvalues.HasWeightedBoundaryH2.integrable_weight_normSq,
  EulerDipoleEigenvalues.integrable_normSq_of_memLp,
  EulerDipoleEigenvalues.integrable_im_mul_conj_of_memLp}
\leanok
For a weighted boundary triple $(F,D,E)$ and $\lambda\in\R$, the functions
$F$, $D$, $wE$, $A$ and $wE-2\xi D-i\lambda D$ are in $L^2(-1,1)$. Hence
the three integrands of $\mathcal B_\lambda(F,D,E)$ and $|A|^2$, $|D|^2$
are integrable on $(-1,1)$.
\end{lemma}

\begin{proof}
\leanok
$L^2(-1,1)$ is closed under sums and under multiplication by bounded
continuous coefficients, and the product of two $L^2$ functions is
integrable.
\end{proof}

\begin{lemma}[Completion of squares on the weighted domain]
\label{lem:weighted-completion}
\uses{def:boundary-form-xi, def:h1-interval, def:weighted-boundary-domain,
  lem:h1-interval-endpoints}
\lean{EulerDipoleEigenvalues.HasWeightedBoundaryH2.completion,
  EulerDipoleEigenvalues.HasWeightedBoundaryH2.sum_of_squares}
\leanok
Let $(F,D,E)$ be a weighted boundary triple with traces $\widetilde F(\pm1)$
(Lemma~\ref{lem:h1-interval-endpoints}), and $\lambda\in\R$. With
$A=wE-\xi D+F-i\lambda D$,
\[
  \mathcal B_\lambda(F,D,E)
  =\frac12\int_{-1}^1|A|^2\,d\xi
   +\frac12\int_{-1}^1|D|^2\,d\xi
   -\frac12\bigl(|\widetilde F(1)|^2+|\widetilde F(-1)|^2\bigr).
\]
In particular, for zero traces $\mathcal B_\lambda(F,D,E)$ is a sum of squares.
\end{lemma}

\begin{proof}
\uses{lem:boundary-flux, lem:weighted-boundary-integrable,
  lem:weighted-boundary-local, lem:l2-endpoint-liminf}
\leanok
Let $-1<c<d<1$. On $[c,d]$ both $F$ and $D$ are absolutely continuous, with
$F'=D$ and $D'=E$ almost everywhere
(Lemma~\ref{lem:weighted-boundary-local}). Hence $K$ is absolutely
continuous on $[c,d]$, and by Lemma~\ref{lem:boundary-flux}
\[
  \int_c^d\Bigl(\tfrac12|wE-2\xi D-i\lambda D|^2-w|D|^2
     +\lambda\operatorname{Im}(D\bar F)\Bigr)
  =\frac12\int_c^d|A|^2+\frac12\int_c^d|D|^2+K(d)-K(c).
\]
By Lemma~\ref{lem:l2-endpoint-liminf} applied to $D$, choose $d_n\uparrow1$
and $c_n\downarrow-1$ with $w|D|^2\to0$ along both sequences (note
$w\le2(1-\xi)$ and $w\le2(1+\xi)$). Since $F=\widetilde F$ in the interior
and $\widetilde F$ is continuous on $[-1,1]$, $F$ stays bounded, so
$|w\operatorname{Re}(D\bar F)|\le\sqrt w\,(w|D|^2)^{1/2}|F|\to0$ as well.
Therefore $K(d_n)\to-\tfrac12|\widetilde F(1)|^2$ and
$K(c_n)\to\tfrac12|\widetilde F(-1)|^2$. The three integrals converge to
their values on $(-1,1)$ by dominated convergence
(Lemma~\ref{lem:weighted-boundary-integrable}).

In Lean the sequences are not constructed: an exhaustion lemma
(\texttt{intervalIntegral\_eq\_of\_frequently}) shows that if $P$ is a
primitive of an integrable $g$ on $(a,b)$ and $P$ comes arbitrarily close to
$L_a$, $L_b$ near the endpoints, then $\int_a^bg=L_b-L_a$. Indeed, the
primitive $P(m)+\int_m^xg$ converges at each endpoint, and its limit must be
the cluster value.
\end{proof}

\begin{lemma}[Affine splitting on the weighted domain]
\label{lem:weighted-affine-splitting}
\uses{def:boundary-form-xi, def:h1-interval, def:weighted-boundary-domain,
  lem:h1-interval-endpoints}
\lean{EulerDipoleEigenvalues.HasWeightedBoundaryH2.affine_splitting,
  EulerDipoleEigenvalues.HasWeightedBoundaryH2.extension_subAffine,
  EulerDipoleEigenvalues.HasWeightedBoundaryH2.integral_boundarySquare_eq,
  EulerDipoleEigenvalues.boundaryPrimitive,
  EulerDipoleEigenvalues.hasDerivAt_boundaryPrimitive}
\leanok
Let $(F,D,E)$ be a weighted boundary triple, $\lambda\in\R$, and $a,b\in\C$
with $\widetilde F(-1)=a-b$ and $\widetilde F(1)=a+b$. Then
$(F_0,D_0,E)=(F-a-b\xi,D-b,E)$ is a weighted boundary triple with zero
traces, and
\[
  \mathcal B_\lambda(F,D,E)=\mathcal B_\lambda(F_0,D_0,E)
    +\lambda^2|b|^2+2\lambda\operatorname{Im}(b\bar a).
\]
\end{lemma}

\begin{proof}
\uses{lem:weighted-boundary-linear, lem:weighted-completion,
  lem:weighted-boundary-integrable, lem:weighted-boundary-local,
  lem:l2-endpoint-liminf}
\leanok
The triple property is Lemma~\ref{lem:weighted-boundary-linear}. Let
$A_0=wE-\xi D_0+F_0-i\lambda D_0$. Both $\int_{-1}^1D_0$ and
$\int_{-1}^1A_0$ vanish. For the first, it equals
$\widetilde F_0(1)-\widetilde F_0(-1)=0$. For the second, on interior
intervals $Q_0=wD_0+\xi F_0-i\lambda F_0$ is absolutely continuous with
$Q_0'=A_0$ almost everywhere; along the sequences of
Lemma~\ref{lem:l2-endpoint-liminf} for $D_0$, $w|D_0|\to0$ and
$F_0\to0$, so $Q_0\to0$ at both ends.

Since $A=A_0+(a-i\lambda b)$ and $D=D_0+b$, the cross terms vanish and,
by Lemma~\ref{lem:weighted-completion} applied to both triples,
\[
  \mathcal B_\lambda(F,D,E)-\mathcal B_\lambda(F_0,D_0,E)
  =|a-i\lambda b|^2+|b|^2-\tfrac12\bigl(|a+b|^2+|a-b|^2\bigr)
  =\lambda^2|b|^2+2\lambda\operatorname{Im}(b\bar a).
\]
\end{proof}

\begin{lemma}[Positivity on the weighted domain]
\label{lem:weighted-boundary-positivity}
\uses{def:boundary-form-xi, def:h1-interval, def:weighted-boundary-domain,
  lem:h1-interval-endpoints}
\lean{EulerDipoleEigenvalues.HasWeightedBoundaryH2.boundaryForm_nonneg}
\leanok
For a weighted boundary triple $(F,D,E)$, real $\lambda$, and
$\widetilde F(1)=\pm\widetilde F(-1)$, one has
$\mathcal B_\lambda(F,D,E)\ge0$.
\end{lemma}

\begin{proof}
\uses{lem:weighted-affine-splitting, lem:weighted-completion}
\leanok
Take the affine part with coefficients
$a=(\widetilde F(1)+\widetilde F(-1))/2$ and
$b=(\widetilde F(1)-\widetilde F(-1))/2$. The zero-trace remainder has
$\mathcal B_\lambda(F_0,D_0,E)=\frac12\int|A_0|^2+\frac12\int|D_0|^2\ge0$
by Lemma~\ref{lem:weighted-completion}, and the cross term
$2\lambda\operatorname{Im}(b\bar a)$ of
Lemma~\ref{lem:weighted-affine-splitting} vanishes since $a=0$ or $b=0$.
\end{proof}

\begin{remark}[Digest note: no density argument is needed]
The manuscript extends the smooth sum of squares to $\mathcal D_w$ ``by the
same density argument''. Lemma~\ref{lem:weighted-completion} instead
exhausts the interval and controls the endpoint flux directly. No smooth
approximation or Hardy inequality is needed, and the identity includes
the endpoint correction for arbitrary value traces.
\end{remark}

\subsection{Boundary form consequences}

\begin{lemma}[Sum of squares, extension]
\label{lem:trace-sos-extension}
\uses{lem:trace-sos-smooth, def:h1-interval, def:h1-zero-interval,
 def:weighted-boundary-domain, lem:h1-interval-endpoints}
\lean{EulerDipoleEigenvalues.HasWeightedBoundaryH2.sum_of_squares}
\leanok
The identity of Lemma~\ref{lem:trace-sos-smooth} holds for $F\in\mathcal D_w$
with zero endpoint traces (equivalently, $F\in H^1_0(-1,1)$ with the
weighted second-order regularity of Definition~\ref{def:weighted-boundary-domain}).
\end{lemma}

\begin{proof}
\uses{lem:boundary-form-xi, lem:weighted-completion}
\leanok
By Lemma~\ref{lem:boundary-form-xi}, $\mathcal F_\lambda(F)=\mathcal B_\lambda(F,D,E)$.
Apply Lemma~\ref{lem:weighted-completion}; the endpoint correction vanishes.
The manuscript argues by density instead; see the digest note after
Lemma~\ref{lem:weighted-boundary-positivity}.
\end{proof}

\subsection{The canonical unknown for even--even eigenfunctions}

\begin{lemma}[Absolute continuity of the arc traces]
\label{lem:arc-trace-ac}
\uses{def:separatrix-arc, def:stream-lifts}
\lean{EulerDipoleEigenvalues.absolutelyContinuous_contStream_separatrixArc,
  EulerDipoleEigenvalues.AbsolutelyContinuousOnInterval.comp_antitoneOn,
  EulerDipoleEigenvalues.AbsolutelyContinuousOnInterval.comp_monotoneOn}
\leanok
For every $\omega\in L^2$, the trace $\tilde\psi\circ\gamma$ of the continuous stream function
on the separatrix arc is absolutely continuous on every compact interval and differentiable
almost everywhere.
\end{lemma}

\begin{proof}
\uses{lem:trace-h2, lem:lift-sobolev-local, lem:elliptic-torus, lem:lift-representatives,
  lem:separatrix-arc}
\leanok
In diagonal coordinates the arc is the segment $b=\pi/2$, $a=\alpha(s)-\pi/2$ with
$\alpha(s)=2\arctan e^{-s}$. By Lemmas~\ref{lem:elliptic-torus} and~\ref{lem:lift-sobolev-local},
$\tilde\psi\in H^2$ on a bounded neighbourhood of the line (changing $\psi$ on a null set to
$\tilde\psi$ does not affect this), so by Lemma~\ref{lem:trace-h2} the continuous function
$\tilde\psi$ is absolutely continuous along the line. The change of variable $s\mapsto a$ is
decreasing and $1$-Lipschitz, which preserves absolute continuity: disjoint intervals go to
disjoint intervals of at most the same total length. Absolutely continuous functions have
bounded variation and are therefore differentiable almost everywhere.
\end{proof}

\begin{lemma}[Arc regularity of an even torus $H^2$ function]
\label{lem:even-h2-arc-regularity}
\uses{def:sobolev-torus, def:sectors, def:continuous-representative,
 def:cell-arcs, def:torus-of-lift}
\lean{EulerDipoleEigenvalues.even_sobolev_arc_regularity}
\leanok
Let $\phi\in H^2(\T^2)$ be even under both reflections. On each cell arc,
$g_k(t)=\widetilde\phi(\pi(\gamma_k(t)))$ is absolutely continuous on
compact intervals, differentiable almost everywhere, and satisfies
$\int_\R\cosh^2t\,|g_k'(t)|^2dt<\infty$.
No eigenvalue equation is assumed.
\end{lemma}

\begin{proof}
\uses{lem:h2-laplacian-reconstruction, lem:laplacian-symmetries,
 lem:arc-trace-ac, lem:ee-arc-energy}
\leanok
Put $\zeta=\Delta\phi\in L^2$. It is even under both reflections, and
$\widetilde\phi=\widetilde{\Delta^{-1}\zeta}+\hat\phi_0$.
Apply the stream arc estimates to $\zeta$; adding a constant preserves
absolute continuity and does not change the arc derivative.
\end{proof}

\begin{lemma}[Canonical profiles in the weighted interval domain]
\label{lem:canonical-boundary-regularity}
\uses{def:canonical-profile, def:h1-interval, def:weighted-boundary-domain,
 lem:h1-interval-endpoints}
\lean{EulerDipoleEigenvalues.canonical_boundary_regularity}
\leanok
Let $\lambda\in\R\setminus\{0\}$ and let $g:\R\to\C$ be bounded,
locally absolutely continuous and differentiable almost everywhere.
Assume $g'\in L^1$, $\int\cosh^2t\,|g'|^2dt<\infty$,
$\int g'(t)e^{-i\lambda t}dt=0$, and $g(\pm\infty)=c_\pm\in\C$.
For the canonical unknown $z$, put $F=z\circ\operatorname{artanh}$.
Then $F\in\mathcal D_w$ with derivatives $D,E$,
$G=-iwD-\lambda F\in H^1(-1,1)$, and
\[
 G(x)=g(\operatorname{artanh}x)\quad(-1<x<1),\qquad
 F(\pm1)=-c_\pm/\lambda
\]
where endpoint values are $H^1$ traces.
\end{lemma}

\begin{proof}
\uses{lem:canonical-profile-ac, lem:canonical-profile-ode,
 lem:canonical-profile-weighted}
\leanok
The profile $\Omega$ is bounded by $\|g'\|_{L^1}$, locally absolutely
continuous, and tends to zero at both ends. Its weighted bound and
$\cosh t\le e^{|t|}\le2\cosh t$ give $\int\cosh^2t\,|\Omega|^2dt<\infty$.
At $t=\operatorname{artanh}x$, $w=1-x^2$, take
\[
 D=-i\Omega/w,\qquad E=-i\big((i\lambda+2x)\Omega-g'\big)/w^2.
\]
The profile equations and chain rule give $F'=D$, $D'=E$ almost
everywhere, with local absolute continuity on compact subintervals.
Boundedness gives $F\in L^2$; using $dx=dt/\cosh^2t$ gives
$D,wE\in L^2$, since $|i\lambda+2x|\le|\lambda|+2$.
The explicit formulas identify $G=g\circ\operatorname{artanh}$ pointwise
on the open interval, with
$\int|G'|^2dx=\int\cosh^2t\,|g'|^2dt$.
Finally $z(t)\to-c_\pm/\lambda$, and continuity of the $H^1$ extension
identifies its endpoint traces.
\end{proof}

\begin{lemma}[Regularity of the boundary unknown]
\label{lem:ee-boundary-regularity}
\uses{def:boundary-form, def:canonical-profile, def:h1-interval, def:weighted-boundary-domain, lem:h1-interval-endpoints, def:cell-arcs, def:stream-lifts, def:sectors, def:rotation-characters, def:eigen-problem}
\lean{EulerDipoleEigenvalues.ee_boundary_regularity}
\leanok
Let $\omega$ be an even--even eigenfunction with $\lambda\ne0$ and $R^*\omega=\kappa\omega$. On each
arc $\gamma_k$ let $g=\tilde\psi\circ\gamma_k$, $z=-(\Omega+g)/\lambda$ its canonical transport
unknown, and $F(\xi)=z(\operatorname{artanh}\xi)$. Then $F\in\mathcal D_w$ with derivatives
$D,E$, and $G=-iwD-\lambda F\in H^1(-1,1)$. The endpoint values of $F$ are $-c_\pm/\lambda$,
where $c_+=\tilde\psi(R^k(0,-\pi))$ and $c_-=\tilde\psi(R^k(\pi,0))$ are the saddle values at
the ends of $\gamma_k$, and $c_+=\kappa c_-$.
\end{lemma}

\begin{proof}
\uses{lem:ee-arc-energy, lem:ee-transport-traces, lem:arc-trace-ac,
  lem:canonical-profile-ac, lem:canonical-profile-weighted,
  lem:canonical-profile-ode, lem:separatrix-arc,
  lem:continuous-representative, lem:boundary-form-xi}
\leanok
By the rotation character, $g=\kappa^k\,\tilde\psi\circ\gamma$, so $g$ is locally absolutely
continuous and differentiable almost everywhere (Lemma~\ref{lem:arc-trace-ac}), with
$g'\in L^1$, $\int g'e^{-i\lambda t}dt=0$ (Lemma~\ref{lem:ee-transport-traces}) and
$\int\cosh^2s\,|g'|^2ds<\infty$ (Lemma~\ref{lem:ee-arc-energy}). Hence
$\Omega$ is bounded by $\|g'\|_{L^1}$, locally absolutely continuous, satisfies
$\Omega'=i\lambda\Omega-g'$ and $z'=-i\Omega$ almost everywhere, tends to $0$ at $\pm\infty$,
and $\int\cosh^2|\Omega|^2<\infty$ by Lemma~\ref{lem:canonical-profile-weighted}
with $\cosh s\le e^{|s|}\le2\cosh s$.

With $\xi=\tanh s$ and $w=1-\xi^2=\operatorname{sech}^2s$ put
$D=-i\Omega(\operatorname{artanh}\xi)/w$ and
$E=-i\bigl((i\lambda+2\xi)\Omega-g'\bigr)/w^2$ at $s=\operatorname{artanh}\xi$. Then $F'=D$ and
$D'=E$ almost everywhere by the chain rule ($\operatorname{artanh}'=1/w$, $(1/w)'=2\xi/w^2$), and
$d\xi=ds/\cosh^2s$ turns the $L^2$ norms of $F$, $D$ and $wE$ into $\int|z|^2\operatorname{sech}^2$,
$\int\cosh^2|\Omega|^2$ and $\int\cosh^2|(i\lambda+2\xi)\Omega-g'|^2$, all finite. Absolute
continuity on compact subintervals of $(-1,1)$ is preserved by $\operatorname{artanh}$ and by
products with the smooth $1/w$. Since $wD=-i\Omega$, $G=-iwD-\lambda F=g\circ\operatorname{artanh}$,
whose derivative has $\int|G'|^2d\xi=\int\cosh^2|g'|^2ds$. The endpoint values follow from
$z(s)\to-(0+c_\pm)/\lambda$ as $s\to\pm\infty$ and continuity of the $H^1$ extension; finally
$R(\pi,0)=(0,-\pi)$ and $\tilde\psi\circ R=\kappa\tilde\psi$ give $c_+=\kappa c_-$.
\end{proof}

\begin{remark}[Inputs of boundary regularity]
Local absolute continuity of $g$ comes from Lemma~\ref{lem:arc-trace-ac};
the Melnikov identity is supplied by Lemma~\ref{lem:ee-transport-traces}.
The arc energy controls $G=g\circ\operatorname{artanh}$ in $H^1$, while
the weighted profile estimate controls $F'$ and $wF''$. These are
distinct inputs; integrability of $g'$ alone does not replace them.
\end{remark}

\begin{lemma}[Positivity of the boundary form]
\label{lem:ee-trace-positivity}
\uses{def:boundary-form, def:h1-interval, def:weighted-boundary-domain, lem:h1-interval-endpoints}
\lean{EulerDipoleEigenvalues.HasWeightedBoundaryH2.flowBoundaryForm_nonneg}
\leanok
If $F\in\mathcal D_w$ and $F(1)=\pm F(-1)$, then
$\mathcal F_\lambda(F)\ge0$ (manuscript \texttt{(ee:trace-positivity)}).
\end{lemma}

\begin{proof}
\uses{lem:boundary-form-xi, lem:weighted-boundary-positivity}
\leanok
By Lemma~\ref{lem:boundary-form-xi},
$\mathcal F_\lambda(F)=\mathcal B_\lambda(F,D,E)$ for the derivatives $D,E$
of $F$, and Lemma~\ref{lem:weighted-boundary-positivity} applies.
\end{proof}

\section{From transport traces to exclusion}
\label{sec:ee-assembly}

The next interfaces connect the raw eigenfunction lifts, continuous arc
traces and explicit derivative arguments of the boundary form. The final
steps combine the virial, trace and positivity results.

\subsection{Identifying the canonical arc form}

\begin{definition}[Canonical arc form]
\label{def:canonical-arc-form}
\uses{def:canonical-profile}
\lean{EulerDipoleEigenvalues.canonicalArcForm}
\leanok
For $\lambda\in\R$ and $g:\R\to\C$, define
\[
  \mathcal A_\lambda(g)=
  \int_\R\tfrac12\cosh^2s\,|g'(s)|^2\,ds
  +\int_\R\operatorname{Re}(\Omega(s)\overline{g(s)})\,ds,
\]
where $\Omega$ is the canonical profile of $g$. As elsewhere, the definition uses
total integrals; the following lemma supplies the analytic hypotheses.
\end{definition}

\begin{lemma}[Canonical arc form and the explicit boundary form]
\label{lem:canonical-arc-form}
\uses{def:canonical-arc-form, def:boundary-form, def:h1-interval, def:weighted-boundary-domain}
\lean{EulerDipoleEigenvalues.canonicalArcForm_eq_flowBoundaryForm}
\leanok
Let $\lambda\ne0$, $g'\in L^1(\R)$ and
$\cosh^2s\,|g'(s)|^2\in L^1(\R)$. Set $z=-(\Omega+g)/\lambda$ and
$F(\xi)=z(\operatorname{artanh}\xi)$. If $(F,D,E)$ is a weighted boundary triple, then
\[
  \mathcal A_\lambda(g)=
  \mathcal F_\lambda\bigl(F\circ\tanh,\ (wD)\circ\tanh,\
       (w(wE-2\xi D))\circ\tanh\bigr).
\]
The three arguments on the right are the explicit value and derivative functions
of Definition~\ref{def:boundary-form}.
\end{lemma}

\begin{proof}
\uses{lem:canonical-profile-ode, lem:weighted-flow-derivatives}
\leanok
Write $z_0=F\circ\tanh=z$, $z_1=(wD)\circ\tanh$ and
$z_2=(w(wE-2\xi D))\circ\tanh$. The weighted chain rules give
$z'=z_1$ and $z_1'=z_2$ almost everywhere. The profile equation and
$g=-\lambda z-\Omega$ first show that $g$ is differentiable almost everywhere;
no separate absolute-continuity assumption on $g$ is needed. The canonical ODE
then gives $z'=-i\Omega$, so $z_1=-i\Omega$ almost everywhere by uniqueness of
derivatives. Both sides are continuous, hence equal everywhere. Differentiating
gives $-iz_2-\lambda z_1=g'$ almost everywhere, while
$-iz_1=-\Omega$ and $-iz_1-\lambda z=g$ everywhere.
Substitution into the two total integrals proves the identity, with a plus sign
in the second term because $Pz=-\Omega$.
\end{proof}

\subsection{Lift representatives and transport traces}

\begin{lemma}[An invariant neighbourhood of the cell]
\label{lem:cell-neighbourhood}
\uses{def:cells, def:transport-pair}
\lean{EulerDipoleEigenvalues.cellNeighbourhood,
  EulerDipoleEigenvalues.isOrbitUnion_cellNeighbourhood,
  EulerDipoleEigenvalues.isBounded_cellNeighbourhood,
  EulerDipoleEigenvalues.closure_cellSquare_subset_cellNeighbourhood,
  EulerDipoleEigenvalues.rotLift_mem_cellNeighbourhood}
\leanok
The square $U=\{|a|<3\pi/2,\ |b|<3\pi/2\}$ in diagonal coordinates is a bounded open
union of orbits. It contains the closed cell and is invariant under the rotation
$(x,y)\mapsto(y,-x)$.
\end{lemma}

\begin{proof}
\uses{lem:dipole-flow, lem:cell-square}
\leanok
In diagonal coordinates $\dot a=-\cos a\sin b$. If $\cos c=0$, the constant $c$ solves
this scalar equation, whose right side is $1$-Lipschitz in $a$, so by uniqueness an orbit
with $a\ne c$ at time $0$ never reaches $a=c$. The same holds for $b$, since
$\dot b=\sin a\cos b$. Hence the lines $a=\pm3\pi/2$, $b=\pm3\pi/2$ are barriers and, by
continuity, an orbit starting in $U$ stays in $U$. The rotation swaps $a$ and $-b$.
\end{proof}

\begin{lemma}[Representatives of the eigenfunction lifts]
\label{lem:lift-representatives}
\uses{def:stream-lifts, def:rotation-characters}
\lean{EulerDipoleEigenvalues.measurePreserving_torusProj_square,
  EulerDipoleEigenvalues.integrableOn_norm_liftL2_sq,
  EulerDipoleEigenvalues.ae_torusProj,
  EulerDipoleEigenvalues.streamLift_ae_eq_contStream,
  EulerDipoleEigenvalues.continuous_contStream,
  EulerDipoleEigenvalues.contStream_rotLift,
  EulerDipoleEigenvalues.liftL2_rotLift_ae}
\leanok
The covering map is measure preserving from each translate of $(0,2\pi]^2$ onto $\T^2$, so
$|g|^2$ is integrable on $(-2\pi,2\pi]^2\supset\overline\Cell$ for $g\in L^2(\T^2)$, and
almost-everywhere statements on $\T^2$ lift to $\R^2$. For $\omega\in L^2$, $\psi=\tilde\psi$
almost everywhere on $\R^2$ and $\tilde\psi$ is continuous. If $\omega$ has rotation
character $\kappa$, then $\tilde\psi(y,-x)=\kappa\tilde\psi(x,y)$ for every $(x,y)$, and
$\omega(y,-x)=\kappa\omega(x,y)$ almost everywhere.
\end{lemma}

\begin{proof}
\uses{lem:continuous-representative, lem:symmetries-commute}
\leanok
The covering map is measure preserving on each translate of a fundamental square, and
$\R^2$ is a countable union of these. Rotation commutes with $\Delta^{-1}$, so $\psi$ has
the rotation character of $\omega$ almost everywhere. Two continuous functions that agree
almost everywhere on $\T^2$ agree everywhere, since Haar measure charges open sets.
\end{proof}

\begin{lemma}[Symmetries of the continuous stream function]
\label{lem:continuous-stream-symmetries}
\uses{def:stream-lifts, def:sectors}
\lean{EulerDipoleEigenvalues.continuousRep_invLap_comp,
  EulerDipoleEigenvalues.contStream_rotate, EulerDipoleEigenvalues.contStream_iterate_rotate,
  EulerDipoleEigenvalues.contStream_reflect₁, EulerDipoleEigenvalues.contStream_add_two_pi,
  EulerDipoleEigenvalues.inSector_iterate_rotate}
\leanok
For every $\omega\in L^2$ and every $(x,y)$, the continuous stream function of the rotated
vorticity $\omega\circ R$ is $\tilde\psi\circ R$, and $\tilde\psi$ is $2\pi$-periodic in each
variable. If $\omega$ is even under $S_1$, then $\tilde\psi(\pi-y,\pi-x)=\tilde\psi(x,y)$ for every
$(x,y)$. If $\omega$ is even under $S_1$ and $S_2$, so is $\omega\circ R^k$.
\end{lemma}

\begin{proof}
\uses{lem:continuous-representative, lem:symmetries-commute}
\leanok
If $\Delta^{-1}\omega\circ T=\Delta^{-1}\omega'$ almost everywhere for a continuous
measure-preserving $T$, the continuous representatives agree almost everywhere and hence
everywhere. $R$ and $S_1$ commute with $\Delta^{-1}$; $R$ exchanges the two reflections.
\end{proof}

\begin{lemma}[Rotation of the normal and tangential derivatives]
\label{lem:normal-tangent-rotation}
\uses{def:normal-tangent, def:stream-derivatives}
\lean{EulerDipoleEigenvalues.measurePreserving_rotLift, EulerDipoleEigenvalues.liftL2_rotate_ae,
  EulerDipoleEigenvalues.streamPartial_rotate, EulerDipoleEigenvalues.streamNormal_rotate_ae,
  EulerDipoleEigenvalues.streamNormal_iterate_rotate_ae}
\leanok
For every $\omega\in L^2$ and $k\ge0$, the normal and tangential derivatives of the stream
function of $\omega\circ R^k$ are those of $\psi$ composed with $R^k$, almost everywhere.
\end{lemma}

\begin{proof}
\uses{lem:stream-derivative-parities}
\leanok
$\partial_x(\psi\circ R)=-\psi_y\circ R$ and $\partial_y(\psi\circ R)=\psi_x\circ R$ (on Fourier
coefficients), $\rho\circ R=\rho$, and $R$ maps $(\sin x,\sin y)$ to $(\sin y,-\sin x)$; so the
components of $\nabla(\psi\circ R)$ in the frame $(n,Y/\rho)$ at $p$ are those of $\nabla\psi$ at
$Rp$. Induct on $k$; $R$ is measure preserving on $\R^2$.
\end{proof}

\begin{lemma}[Transport pairs with the continuous stream function]
\label{lem:transport-pair-continuous}
\uses{def:transport-pair, def:stream-lifts, def:weak-derivative-local, def:sobolev-local}
\lean{EulerDipoleEigenvalues.HasWeakPartialOn.congr_value,
  EulerDipoleEigenvalues.MemSobolevOn.congr_value,
  EulerDipoleEigenvalues.HasWeakL0On.congr_value,
  EulerDipoleEigenvalues.isTransportPair_contStream}
\leanok
Weak derivatives, local Sobolev membership and weak transport are unchanged by an
almost-everywhere change of representative. Hence for an eigenfunction with $\lambda\ne0$
and a bounded open union of orbits $U$, $(\omega,\tilde\psi)$ is a transport pair on $U$.
\end{lemma}

\begin{proof}
\uses{lem:transport-pair-eigen, lem:lift-representatives}
\leanok
All three notions are defined by integrals against test functions. Apply this to the pair
$(\omega,\psi)$ of Lemma~\ref{lem:transport-pair-eigen} and $\psi=\tilde\psi$ almost
everywhere.
\end{proof}

\begin{lemma}[Eigenfunction transport traces]
\label{lem:ee-transport-traces}
\uses{def:eigen-problem, def:sectors, def:rotation-characters, def:stream-lifts,
  def:cell-arcs, def:level-integral, def:canonical-profile}
\lean{EulerDipoleEigenvalues.ee_transport_traces, EulerDipoleEigenvalues.ee_melnikov,
  EulerDipoleEigenvalues.contStream_iterate_rotLift}
\leanok
Let $\omega$ be an even--even eigenfunction with $\lambda\ne0$ and rotation character
$\kappa$. Put $g_k=\tilde\psi\circ\gamma_k$. Then $g_k'\in L^1(\R)$ and
$\int_\R g_k'(t)e^{-i\lambda t}\,dt=0$ (the Melnikov identity) for every $k$.
For every null set $N$ there are $h_j\in(0,1)\setminus N$, $h_j\to0$, such that
\[
  \oint_{\{\Psi=2h_j\}}\operatorname{Re}(\omega\bar\psi)\,dt
   \longrightarrow\sum_{k=0}^3
      \int_\R\operatorname{Re}(\Omega_k\bar g_k)\,ds .
\]
The level pairing uses the raw $L^2$ representatives; the arc traces use the continuous
stream representative.
\end{lemma}

\begin{proof}
\uses{lem:cell-neighbourhood, lem:transport-pair-continuous, lem:lift-representatives,
  lem:melnikov, lem:pairing-limit, lem:level-integral-congr}
\leanok
Let $U$ be the neighbourhood of Lemma~\ref{lem:cell-neighbourhood}; by
Lemma~\ref{lem:transport-pair-continuous}, $(\omega,\tilde\psi)$ is a transport pair on $U$.
By Lemma~\ref{lem:lift-representatives}, $\tilde\psi$ has the rotation character
$\kappa$ pointwise and $\omega$ almost everywhere, which are the hypotheses of
Lemmas~\ref{lem:melnikov} and~\ref{lem:pairing-limit}. Melnikov gives
$(\tilde\psi\circ\gamma)'\in L^1$ on the separatrix arc $\gamma$, and
$\tilde\psi\circ\gamma_k=\kappa^k\,\tilde\psi\circ\gamma$ transfers this to all four arcs.
The pairing limit is stated with $\tilde\psi$ in the level integrals, while the
lemma uses the raw $\psi$. By Lemma~\ref{lem:level-integral-congr} the two level
integrals agree outside a null set $N'$ of levels; apply the pairing limit with the
null set $N\cup N'$.
\end{proof}

\begin{remark}[Scope of the transport interface]
The proof of Lemma~\ref{lem:ee-transport-traces} uses the rotation character
$\kappa=\pm1$, not the even--even reflection parity. The explicit invariant
neighbourhood avoids a flow-saturation construction. Changes between raw
level representatives and continuous arc traces are handled separately
from the identification of the canonical and explicit boundary forms.
\end{remark}

\subsection{Positivity and cell vanishing}

\begin{lemma}[Positivity of the canonical arc contribution]
\label{lem:ee-canonical-arc-nonneg}
\uses{def:canonical-arc-form, def:eigen-problem, def:sectors, def:rotation-characters,
  def:stream-lifts, def:cell-arcs}
\lean{EulerDipoleEigenvalues.ee_canonicalArcForm_nonneg}
\leanok
Under the eigenfunction and character hypotheses above, if $g_k'\in L^1$ and
$\cosh^2s\,|g_k'|^2\in L^1$, then $\mathcal A_\lambda(g_k)\ge0$.
\end{lemma}

\begin{proof}
\uses{lem:ee-boundary-regularity, lem:canonical-arc-form, lem:ee-trace-positivity}
\leanok
Boundary regularity gives a weighted triple for the canonical unknown with equal or
opposite endpoint traces. Its flow-time form is nonnegative. Identify that form with
$\mathcal A_\lambda(g_k)$ by Lemma~\ref{lem:canonical-arc-form}.
\end{proof}

\begin{lemma}[Null cosine diagonal]
\label{lem:cos-diagonal-null}
\lean{EulerDipoleEigenvalues.ae_cos_fst_ne_cos_snd}
\leanok
The set $\{(x,y)\in\R^2:\cos x=\cos y\}$ is null.
\end{lemma}

\begin{proof}
\leanok
For each $x$ its vertical section is the countable set
$\{2m\pi+x,2m\pi-x:m\in\Z\}$. Apply Fubini to the measurable set.
\end{proof}

\begin{lemma}[Nonnegativity of the evaluated bulk]
\label{lem:ee-bulk-density-nonneg}
\uses{def:ee-bulk-density, def:cells}
\lean{EulerDipoleEigenvalues.eeBulkDensity_nonneg}
\leanok
For every $\omega\in L^2$, $\mathcal J_\omega\ge0$ on the entire lifted cell,
including its center.
\end{lemma}

\begin{proof}
\uses{lem:normal-field-algebra, lem:ee-density, lem:dipole-critical-points, lem:cell-square}
\leanok
Away from the center use Lemma~\ref{lem:ee-density}. At $\Psi=2$, the total definitions
give $f=c=0$, and the remaining term is $2|\psi_n|^2\ge0$.
\end{proof}

\begin{lemma}[Vorticity from zero bulk density]
\label{lem:ee-vorticity-from-density}
\uses{def:ee-bulk-density, def:normal-field, def:cells, def:torus-of-lift}
\lean{EulerDipoleEigenvalues.normalCoeff_pos_of_cos_ne,
  EulerDipoleEigenvalues.lift_ae_zero_of_eeBulkDensity_ae_zero}
\leanok
Where $\Psi>0$ and $\cos x\ne\cos y$, the coefficient $f$ is strictly positive.
Consequently, if $\mathcal J_\omega=0$ almost everywhere on the lifted cell, then
the lift of $\omega$ is zero almost everywhere there.
\end{lemma}

\begin{proof}
\uses{lem:normal-field-algebra, lem:cos-diagonal-null, lem:cell-square}
\leanok
The cosine inequality excludes $\Psi=2$, so $0<\Psi<2$ and $\rho^2>0$.
The formula $f=\Psi(2-\Psi)(\cos x-\cos y)^2/\rho^4$ is then strictly positive.
All terms in $\mathcal J_\omega=f|\omega|^2+(2-c)|\psi_n|^2+c|\psi_t|^2$
are nonnegative. Thus zero density forces $\omega=0$ off the null cosine diagonal.
\end{proof}

\begin{lemma}[Vanishing on the lifted positive cell]
\label{lem:ee-cell-vanishing}
\uses{def:sectors, def:cells, def:torus-of-lift}
\lean{EulerDipoleEigenvalues.InSector.eq_zero_of_lift_ae_zero_on_cellSquare}
\leanok
If $\omega$ lies in any reflection sector and its lift vanishes almost
everywhere on the lifted positive cell, then $\omega=0$ in $L^2(\T^2)$.
No eigenvalue or stream-function relation is needed.
\end{lemma}

\begin{proof}
\uses{lem:cell-square, lem:reflection-algebra, lem:separatrix-lines}
\leanok
The projection from the fundamental square $(-\pi,\pi]^2$ preserves measure and is
injective. Its positive part is exactly the lifted positive cell, so $\omega=0$
almost everywhere on the positive torus cell. The separatrix is null: at each fixed
$x$ its section has at most two points. Finally $S_1$ preserves measure, exchanges
the two cells, and acts by multiplication by $\pm1$ on $\omega$, giving
vanishing on the negative cell.
\end{proof}

\begin{lemma}[Nonnegative exhaustion without global integrability]
\label{lem:nonnegative-exhaustion}
\lean{EulerDipoleEigenvalues.ae_zero_of_truncated_integral_tendsto}
\leanok
Let $(X,\mu)$ be a measure space, $s$ measurable, $r:X\to\R$ measurable and positive on
$s$, and $J:X\to\R$ nonnegative on $s$. Suppose $t_j>0$, $t_j\to0$, and $J$ is integrable
on each $s_j=\{x\in s:r(x)>t_j\}$. If $\int_{s_j}J\to b\le0$, then $J=0$ almost everywhere
on $s$.
\end{lemma}

\begin{proof}
\leanok
For fixed $j$, eventually $t_m<t_j$, hence $s_j\subset s_m$ and
$0\le\int_{s_j}J\le\int_{s_m}J$. Passing to the limit gives $\int_{s_j}J=0$, hence
$J=0$ almost everywhere on $s_j$. The countable union of the $s_j$ is $s$.
Neither monotonicity of the sequence nor integrability on all of $s$ is needed.
\end{proof}

\begin{lemma}[Even--even exclusion in one rotation character]
\label{lem:ee-character-exclusion}
\uses{def:eigen-problem, def:sectors, def:rotation-characters}
\lean{EulerDipoleEigenvalues.even_even_character_exclusion}
\leanok
An even--even eigenfunction with $\lambda\ne0$ and rotation character $\kappa=\pm1$
vanishes.
\end{lemma}

\begin{proof}
\uses{lem:ee-truncated-virial, lem:ee-gradient-traces, lem:ee-transport-traces,
  lem:ee-canonical-arc-nonneg, lem:ee-bulk-density-nonneg, lem:nonnegative-exhaustion,
  lem:ee-vorticity-from-density, lem:ee-cell-vanishing}
\leanok
Let $N$ be the union of the exceptional sets for the gradient traces and truncated virial
identity, including truncated integrability. Apply Lemma~\ref{lem:ee-transport-traces}
outside $N$. On the resulting levels, the right-hand side of the truncated identity
converges to $-2\sum_k\mathcal A_\lambda(g_k)\le0$, by the gradient limits and
Lemma~\ref{lem:ee-canonical-arc-nonneg}. Apply Lemma~\ref{lem:nonnegative-exhaustion}
with $r=\Psi$ and $t_j=2h_j$ to conclude that $\mathcal J_\omega=0$ almost everywhere on
the cell. Lemma~\ref{lem:ee-vorticity-from-density} gives vanishing of the vorticity
lift, and Lemma~\ref{lem:ee-cell-vanishing} gives $\omega=0$.
\end{proof}

\section{Conclusion}

\begin{theorem}[Even--even exclusion]
\label{thm:ee-exclusion}
\uses{def:eigen-problem, def:sectors}
\lean{EulerDipoleEigenvalues.even_even_exclusion}
\leanok
If $\omega$ solves the eigenvalue problem with $\lambda\ne0$ and
$S_1^*\omega=S_2^*\omega=\omega$, then $\omega=0$ (manuscript
\texttt{(ee:theorem)}).
\end{theorem}

\begin{proof}
\uses{lem:rotation-characters, lem:ee-character-exclusion}
\leanok
Decompose $\omega=\omega_++\omega_-$ into rotation characters by
Lemma~\ref{lem:rotation-characters}. Both components remain even--even eigenfunctions
at the same $\lambda$. Apply Lemma~\ref{lem:ee-character-exclusion} to each component
and add the resulting zero functions.
\end{proof}

\begin{remark}[Digest note: direct vorticity endgame]
The positive vorticity term gives $\omega=0$ directly from
$\mathcal J_\omega=0$: the coefficient $f$ vanishes only on the null
cosine diagonal in the open cell. There is no need first to show that
$\psi$ is constant or use its mean-zero normalization. Truncated
integrability suffices by Lemma~\ref{lem:nonnegative-exhaustion}.
\end{remark}

\begin{remark}[Digest note]
The manuscript's overview in the introduction normalizes the boundary
unknown by $F(\pm1)=0$ (``suppose also that $\psi$ vanishes at the
saddles'') and handles the general case with the affine splitting. The
$R$-character enters only to guarantee $c_+=\pm c_-$, which kills the cross
term $2\lambda\operatorname{Im}(b\bar a)$. It would be worth checking whether
the $R$-decomposition is needed anywhere else in this sector (it is used in
the Melnikov argument, Lemma~\ref{lem:melnikov}).
\end{remark}
